% Intended LaTeX compiler: pdflatex
\documentclass[a4paper,12pt]{article}
\usepackage[utf8]{inputenc}
\usepackage[T1]{fontenc}
\usepackage{graphicx}
\usepackage{longtable}
\usepackage{wrapfig}
\usepackage{rotating}
\usepackage[normalem]{ulem}
\usepackage{amsmath}
\usepackage{amssymb}
\usepackage{capt-of}
\usepackage{hyperref}
\usepackage{\string~/.emacs.d/common}
\author{mahmood}
\date{\today}
\title{formal proof}
\hypersetup{
 pdfauthor={mahmood},
 pdftitle={formal proof},
 pdfkeywords={},
 pdfsubject={},
 pdfcreator={Emacs 28.2 (Org mode 9.5.5)}, 
 pdflang={English}}
\begin{document}

\maketitle
\begin{note}
this document and others can be found on my website\\
\url{https://mahmoodsheikh36.github.io/}
\end{note}
\begin{definition}
a \textbf{formal proof} is a finite \href{20220711175112-sequence.org}{sequence} of \href{20230524212241-premise.org}{premise}s each of which is an axiom, an assumption, or follows from the preceding statements in the sequence by a \href{20230524202636-rule_of_inference.org}{rule of inference}, which are used to prove \href{20230525215134-argument.org}{argument}s.\\
a proof sequence from \(\Gamma\) to a \href{20230401175139-statement.org}{proposition} \(a\) is one such that \(a\) is the last proposition in the sequence which we call the \href{20230524212250-conclusion.org}{conclusion}. we write \(\Gamma \vdash \alpha\)\\
\end{definition}
usually, on the ride side of each formula we write a description of how it was derived and from which previous lines in the inference process\\
\section{correct use of equivalences and implications}
\label{sec:orgd48ac7b}
use of \href{20230525090415-equivalence.org}{equivalence}s can also be done on subformulas\\
use of \href{20230525090605-logical_implication.org}{logical implication}s cannot be done on subformulas, only on full rows during the formal proof\\
\begin{my_example}
an example of incorrect usage:\\
\begin{align*}
  1.\ & x & P\\
  2.\ & (x \land y) \to z & P\\
  3.\ & x \to z & 2,I3 \quad \text{incorrect}\\
  4.\ & z & 1,3,I11
\end{align*}
\end{my_example}
\end{document}